%%%%%%%%%%%%%%%%%%%%%%%%%%%%%%%%%%%%%%%%%%%%%%%%%%%%%%%%%%%%%%%%%%%
%                                                                 %
%                          CHAPTER FIVE                          %
%                                                                 %
%%%%%%%%%%%%%%%%%%%%%%%%%%%%%%%%%%%%%%%%%%%%%%%%%%%%%%%%%%%%%%%%%%%

\chapter{Πειράματα και Αποτελέσματα}
\label{ch:experiments}

Το παρόν κεφάλαιο παρουσιάζει τη συστηματική πειραματική αξιολόγηση των μεθόδων που 
διατυπώθηκαν στο Κεφάλαιο~\ref{ch:method}. Τα πειράματα οργανώνονται σε ένδεκα ομάδες που ακολουθούν 
μια λογική αλληλουχία. Πρώτα επαληθεύεται η ανταγωνιστικότητα της προσαρμοσμένης 
υλοποίησης έναντι υφιστάμενων βιβλιοθηκών RL. Στη συνέχεια απομονώνεται διαδοχικά η 
επίδραση της αρχιτεκτονικής, της συνάρτησης ανταμοιβής, του συνόλου χαρακτηριστικών, 
του κόστους συναλλαγής, της αντιστοίχισης ενεργειών και του μήκους επεισοδίου. 
Έπειτα εξετάζονται οι περιορισμοί ρευστότητας, τόσο σε μορφή σκληρής ποινής όσο και σε 
Λαγκρανζιανή διατύπωση. Τέλος, η βελτιστοποίηση μέσω CPCV και η ενσωμάτωση σήματος 
συναισθήματος από το εξειδικευμένο LLM αναδεικνύουν τον συνδυαστικό μηχανισμό που οδηγεί στην καλύτερη 
συνολική απόδοση. Η χρήση του σήματος συναισθήματος εξετάζεται με δύο τρόπους, ως χαρακτηριστικό του χώρου καταστάσεων 
και ως χαρακτηριστικό του χώρου καταστάσεων αλλά με επίδραση και στις ενέργειες, όπως αναλύθηκε στην Ενότητα~\ref{}. 
Όλες οι συγκρίσεις βασίζονται στην κοινή μεθοδολογία αξιολόγησης που 
περιγράφεται στην Ενότητα~\ref{sec:exp-validation}.

\section{Επισκόπηση και Μεθοδολογία Αξιολόγησης}
\label{sec:exp-overview}

\subsection{Πειραματικές Ομάδες}
\label{subsec:exp-groups}

Τα πειράματα οργανώνονται σε ένδεκα ομάδες, καθεμία από τις οποίες απομονώνει έναν 
παράγοντα ενώ διατηρεί σταθερές τις υπόλοιπες συνθήκες. Ο Πίνακας~\ref{tab:exp-groups} 
συνοψίζει τις ομάδες, τον ελεγχόμενο παράγοντα και τον κεντρικό ερευνητικό σκοπό κάθε ομάδας.

\begin{table}[htbp]
  \centering
  \caption{Επισκόπηση ένδεκα πειραματικών ομάδων, ελεγχόμενων παραγόντων και ερευνητικών σκοπών.}
  \label{tab:exp-groups}
  \small
  \begin{tabular}{@{}llp{6.5cm}@{}}
    \toprule
    \textbf{Ομάδα Πειράματος} & \textbf{Ελεγχόμενος Παράγοντας} & \textbf{Ερευνητικός Σκοπός} \\
    \midrule
    Σύγκριση βιβλιοθηκών RL   & Βιβλιοθήκη υλοποίησης           & Επαλήθευση ανταγωνιστικότητας προσαρμοσμένης υλοποίησης, εντοπισμός υπερπροσαρμογής βιβλιοθηκών \\
    Σύγκριση αρχιτεκτονικής   & MLP έναντι LSTM, μέγεθος LSTM   & Απομόνωση επίδρασης αναδρομικής αρχιτεκτονικής \\
    Σχεδιασμός ανταμοιβής     & Συνάρτηση ανταμοιβής             & Αποτροπή παθητικής κατάρρευσης μέσω κατάλληλης ανταμοιβής \\
    Σύνολο χαρακτηριστικών    & Preset χαρακτηριστικών (A--E)   & Επίδραση εμπλουτισμένων τεχνικών δεικτών στη γενίκευση \\
    Κόστος συναλλαγής         & Ποσοστό προμήθειας               & Ευαισθησία βέλτιστου πράκτορα σε ρεαλιστικές προμήθειες \\
    Δράσεις και έκθεση        & Αντιστοίχιση ενεργειών           & Επίδραση μηχανισμού αντιστοίχισης στη συμπεριφορά \\
    Μήκος επεισοδίου          & Βήματα ανά επεισόδιο             & Ανταλλαγή μεταξύ χρονικού ορίζοντα και ποικιλομορφίας δειγμάτων \\
    Σκληροί περιορισμοί ρευστ. & Τύπος ποινής ρευστότητας         & Επίδραση σκληρών ορίων ρευστότητας στο προφίλ κινδύνου \\
    Λαγκρανζιανοί περιορισμοί & Διατύπωση Λαγκρανζιανού CMDP    & Αποτελεσματικότητα και παθολογίες Λαγκρανζιανής χαλάρωσης \\
    CPCV βελτιστοποίηση        & Υπερπαράμετροι μέσω CPCV/Optuna & Αξία βελτιστοποίησης υπερπαραμέτρων με χρονική εκκαθάριση \\
    Συναίσθημα FinGPT          & Σήμα LLM συναισθήματος           & Συνέργεια σήματος συναισθήματος με Λαγκρανζιανό περιορισμό \\
    \bottomrule
  \end{tabular}
\end{table}

\subsection{Μεθοδολογία Αξιολόγησης}
\label{sec:exp-validation}

Η αξιολόγηση ακολουθεί τέσσερις αρχές για την αποφυγή υπερπροσαρμογής και την εξασφάλιση 
συγκρισιμότητας αποτελεσμάτων.

\textbf{Χρονικός διαχωρισμός.} Κάθε πράκτορας εκπαιδεύεται σε δεδομένα από 29 Απριλίου 
2019 έως 21 Νοεμβρίου 2024 και αξιολογείται σε αδύνατα δεδομένα από 22 Νοεμβρίου 2024 έως 
9 Νοεμβρίου 2025 (περίοδος συναλλαγών). Αυτός ο ορίζοντας αξιολόγησης δεν επικαλύπτεται με 
κανένα τμήμα εκπαίδευσης.

\textbf{Τυχαιοποίηση.} 
Κάθε πείραμα διεξάγεται πέντε φορές με διαφορετικούς παράγοντες τυχαιοποίησης. 
Ο μέσος όρος και η τυπική απόκλιση των αποδόσεων αυτών χρησιμοποιούνται ως κριτήρια σύγκρισης.

\textbf{Ενιαία αρχική κατάσταση αξιολόγησης.} Κατά την εκπαίδευση, η αρχική κατάσταση 
δειγματοληπτείται τυχαία (βλ.\ Ενότητα~\ref{sec:start_state_effectiveness} του Κεφαλαίου~4), 
εξασφαλίζοντας εργοδικότητα της πολιτικής. Κατά την αξιολόγηση, η αρχική κατάσταση 
καθορίζεται σταθερά ως, το 50\% του κεφαλαίου ($500{,}000\$$) σε μετρητά και 
το υπόλοιπο 50\% ($500{,}000\$$) ισόποσα κατανεμημένο στα 10 περιουσιακά στοιχεία. Αυτό εξαλείφει σημαντική 
πηγή τυχαιότητας στα αποτελέσματα.

\textbf{Κόστη συναλλαγών.} Εφόσον δεν αναφέρεται διαφορετικά, το κόστος 
συναλλαγής ορίζεται σε $10^{-4}$.

\subsection{Βασική Γραμμή Αναφοράς}
\label{subsec:exp-baseline}

Η στρατηγική αναφοράς είναι ισόποση αγορά-και-κράτηση (equal-weight buy-and-hold), δηλαδή 
επένδυση του αρχικού κεφαλαίου ($1{,}000{,}000\$$) ισόποσα στα 10 κρυπτονομίσματα χωρίς 
επανεξισορρόπηση. Για την περίοδο αξιολόγησης επιτυγχάνει Sharpe ratio 0,267, 
τελική αξία $922{,}368\$$, μέγιστη πτώση (Max Drawdown) $-58{,}3\%$ και ετήσια μεταβλητότητα (annualized volatility) 75,4\%. 
Η υψηλή μεταβλητότητα και ο χαμηλός Sharpe αντικατοπτρίζουν τη χαρακτηριστική αστάθεια 
της αγοράς κρυπτονομισμάτων. Αυτές οι τιμές αποτελούν το σημείο αναφοράς σε όλες τις συγκρίσεις.


\section{Σύγκριση Βιβλιοθηκών Ενισχυτικής Μάθησης}
\label{sec:exp-libraries}


\subsection{Κίνητρο και Πειραματικός Σχεδιασμός}
\label{subsec:exp-libraries-design}

Πριν από οποιαδήποτε αρχιτεκτονική ή αλγοριθμική καινοτομία, είναι αναγκαίο να τεθεί ένα 
αξιόπιστο σημείο αναφοράς που να αντικατοπτρίζει την τρέχουσα κατάσταση των εργαλείων. 
Η ομάδα σύγκρισης βιβλιοθηκών αξιολογεί τρεις καθιερωμένες υλοποιήσεις PPO, από ElegantRL, 
RLlib και Stable Baselines3. Και οι τρεις χρησιμοποιούν αρχιτεκτονική PPO-MLP 
με ανταμοιβή PnL και τις προεπιλεγμένες τους υπερπαραμέτρους, εκπαιδευμένες στα ίδια δεδομένα. 
Ο σκοπός είναι διττός: να επαληθευτεί ότι η προσαρμοσμένη υλοποίηση που αναπτύχθηκε στην 
παρούσα εργασία δεν υστερεί των βιβλιοθηκών, και να εντοπιστούν τυχόν παθολογίες που 
συνδέονται με τις προεπιλεγμένες ρυθμίσεις.

\subsection{Αποτελέσματα Περιόδου Αξιολόγησης}
\label{subsec:exp-libraries-test}

Ο Πίνακας~\ref{tab:g2-test} παρουσιάζει τα αποτελέσματα εκτός δείγματος για τους τρεις πράκτορες βιβλιοθηκών και τη βασική γραμμή αγοράς-και-κράτησης.

\begin{table}[htbp]
  \centering
  \caption{Αποτελέσματα εκτός δείγματος για τους πράκτορες βιβλιοθηκών RL (PPO-MLP, ανταμοιβή PnL, περίοδος αξιολόγησης 2024--2025).}
  \label{tab:g2-test}
  \begin{tabular}{@{}lcccc@{}}
    \toprule
    \textbf{Πράκτορας} & \textbf{Sharpe} & \textbf{Μεταβλ.} & \textbf{Max DD} & \textbf{Τελικό (K)} \\
    \midrule
    ElegantRL PPO-MLP     & 0,349 & 0,752 & $-0{,}561$ & 980 \\
    RLlib PPO-MLP         & 0,056 & 0,707 & $-0{,}551$ & 816 \\
    Stable Baselines3 PPO-MLP & 0,222 & 0,695 & $-0{,}530$ & 919 \\
    \midrule
    Buy-and-Hold (αναφορά) & 0,267 & 0,754 & $-0{,}583$ & 922 \\
    \bottomrule
  \end{tabular}
\end{table}

Το Σχήμα~\ref{fig:g2-returns} παρουσιάζει τις σωρευτικές αποδόσεις κατά την περίοδο αξιολόγησης. Ο ElegantRL επιτυγχάνει τη μέγιστη Sharpe ratio μεταξύ των βιβλιοθηκών (0,349) και τελική αξία 980K, υπερβαίνοντας οριακά τη βασική γραμμή. Ο RLlib αποδίδει σημαντικά χειρότερα (Sharpe 0,056, 816K), χαμηλότερα και από τη βασική γραμμή. Η Stable Baselines3 αποδίδει κοντά στη βασική γραμμή (Sharpe 0,222, 919K), χωρίς ουσιαστικό πλεονέκτημα έναντι της παθητικής στρατηγικής.

\begin{figure}[htbp]
  \centering
  \includegraphics[width=0.85\textwidth]{../../out/final_groups/trade/g2/cumulative_returns_mean.png}
  \caption{Σωρευτικές αποδόσεις κατά την περίοδο αξιολόγησης (2024--2025) για τους πράκτορες βιβλιοθηκών RL και τη βασική γραμμή αγοράς-και-κράτησης. Ο ElegantRL υπερτερεί οριακά, ενώ ο RLlib υποαποδίδει σημαντικά.}
  \label{fig:g2-returns}
\end{figure}

\subsection{Παθολογία Υπερπροσαρμογής}
\label{subsec:exp-libraries-overfit}

Η εικόνα αλλάζει ριζικά όταν εξεταστούν τα αποτελέσματα εντός δείγματος, όπως παρουσιάζεται στον Πίνακα~\ref{tab:g2-train}.

\begin{table}[htbp]
  \centering
  \caption{Αποτελέσματα εντός δείγματος για τους πράκτορες βιβλιοθηκών RL (περίοδος εκπαίδευσης 2019--2024). Η τελική αξία για τη Stable Baselines3 αντικατοπτρίζει ακραία υπερπροσαρμογή.}
  \label{tab:g2-train}
  \begin{tabular}{@{}lcc@{}}
    \toprule
    \textbf{Πράκτορας} & \textbf{Sharpe (εκπ.)} & \textbf{Τελική αξία εκπαίδευσης} \\
    \midrule
    ElegantRL PPO-MLP         & 0,981 & 13,3M \\
    RLlib PPO-MLP             & 1,673 & 280,0M \\
    Stable Baselines3 PPO-MLP & 2,948 & $\approx 1{,}37 \times 10^{11}$ \\
    \midrule
    Buy-and-Hold              & 0,868 & 7,99M \\
    \bottomrule
  \end{tabular}
\end{table}

Η Stable Baselines3 επιτυγχάνει Sharpe εκπαίδευσης 2,948 και τελική αξία εντός δείγματος περίπου 137 δισεκατομμύρια νομισματικές μονάδες από αρχικό κεφάλαιο 1.000K, δηλαδή πολλαπλασιασμό κεφαλαίου κατά 137.000x. Αυτή η τιμή είναι εμφανώς αδύνατο να αντιστοιχεί σε πραγματική στρατηγική: αντικατοπτρίζει ακραία υπερπροσαρμογή στα δεδομένα εκπαίδευσης. Το αποτέλεσμα εκτός δείγματος (Sharpe 0,222, 919K) επιβεβαιώνει ότι η in-sample απόδοση δεν έχει γενίκευση. Ο RLlib εμφανίζει παρόμοιο πρότυπο: 280M in-sample αλλά μόλις 816K out-of-sample.

Αυτά τα ευρήματα αποδεικνύουν ότι οι προεπιλεγμένες ρυθμίσεις των βιβλιοθηκών, απουσία CPCV ή άλλου μηχανισμού αποφυγής χρονικής διαρροής, οδηγούν σε σοβαρή υπερπροσαρμογή. Η προσαρμοσμένη υλοποίηση που αναπτύχθηκε στην παρούσα εργασία, όπως θα φανεί στις επόμενες ενότητες, επιτυγχάνει ανώτερη εκτός δείγματος Sharpe (0,676 και άνω) ενώ αποφεύγει αυτά τα ακραία in-sample φαινόμενα.

% ==============================================================
\section{Αρχιτεκτονική Πράκτορα: MLP έναντι LSTM}
\label{sec:exp-arch}
% ==============================================================

\subsection{Κίνητρο και Πειραματικός Σχεδιασμός}
\label{subsec:exp-arch-design}

Η θεωρητική αιτιολόγηση για τη χρήση LSTM στο περιβάλλον διαχείρισης χαρτοφυλακίου ανέφερθηκε στο Κεφάλαιο~3 (\ref{subsec:pomdp}): το πρόβλημα είναι φύσει POMDP, καθώς ο πράκτορας παρατηρεί μόνο τρέχουσες τιμές και τεχνικούς δείκτες, χωρίς άμεση πρόσβαση σε κρυφές μεταβλητές της αγοράς όπως η κατεύθυνση τάσης ή η αστάθεια καθεστώτος. Η LSTM κρυφή κατάσταση $(h_t, c_t)$ λειτουργεί ως προσέγγιση της κατάστασης πεποίθησης.

Για να εκτιμηθεί εμπειρικά αυτό το πλεονέκτημα, η ομάδα αρχιτεκτονικής συγκρίνει τέσσερις διαμορφώσεις: PPO-MLP 2$\times$512 και PPO-LSTM σε τρία μεγέθη κρυφής κατάστασης (2$\times$256, 2$\times$512, 2$\times$1024). Όλες οι διαμορφώσεις χρησιμοποιούν ανταμοιβή PnL και σύνολο χαρακτηριστικών βάσης με ορμή, εξασφαλίζοντας ότι η αρχιτεκτονική είναι ο μόνος μεταβλητός παράγοντας εντός της ομάδας.

\subsection{Αποτελέσματα}
\label{subsec:exp-arch-results}

Ο Πίνακας~\ref{tab:g1-test} παρουσιάζει τα αποτελέσματα εκτός δείγματος.

\begin{table}[htbp]
  \centering
  \caption{Σύγκριση αρχιτεκτονικών MLP και LSTM (ανταμοιβή PnL, χαρακτηριστικά βάσης με ορμή, περίοδος αξιολόγησης 2024--2025).}
  \label{tab:g1-test}
  \begin{tabular}{@{}lcccc@{}}
    \toprule
    \textbf{Αρχιτεκτονική} & \textbf{Sharpe} & \textbf{Μεταβλ.} & \textbf{Max DD} & \textbf{Τελικό (K)} \\
    \midrule
    PPO-MLP 2$\times$512       & 0,313 & 0,782 & $-0{,}559$ & 941 \\
    PPO-LSTM 2$\times$256      & 0,289 & 0,633 & $-0{,}515$ & 982 \\
    \textbf{PPO-LSTM 2$\times$512} & \textbf{0,676} & \textbf{0,687} & $-0{,}467$ & \textbf{1.246} \\
    PPO-LSTM 2$\times$1024     & 0,347 & 0,767 & $-0{,}597$ & 973 \\
    \midrule
    Buy-and-Hold (αναφορά)     & 0,267 & 0,754 & $-0{,}583$ & 922 \\
    \bottomrule
  \end{tabular}
\end{table}

Το Σχήμα~\ref{fig:g1-returns} παρουσιάζει τις σωρευτικές αποδόσεις. Η PPO-LSTM 2$\times$512 επιτυγχάνει Sharpe 0,676 και τελική αξία 1.246K, κατά πολύ ανώτερη τόσο από την PPO-MLP 2$\times$512 (0,313, 941K) όσο και από τη βασική γραμμή (0,267, 922K). Το μέγεθος 2$\times$512 αποδεικνύεται βέλτιστο: το μικρότερο (2$\times$256) δεν έχει επαρκή εκφραστική δύναμη (Sharpe 0,289), ενώ το μεγαλύτερο (2$\times$1024) εμφανίζει χαμηλότερη γενίκευση (Sharpe 0,347), πιθανώς λόγω υπερπροσαρμογής.

\begin{figure}[htbp]
  \centering
  \includegraphics[width=0.85\textwidth]{../../out/final_groups/trade/g1/cumulative_returns_mean.png}
  \caption{Σωρευτικές αποδόσεις εκτός δείγματος για τις αρχιτεκτονικές MLP και LSTM. Η PPO-LSTM 2$\times$512 υπερτερεί σαφώς, ενώ το μεγαλύτερο μέγεθος 2$\times$1024 δεν παρέχει πρόσθετο όφελος.}
  \label{fig:g1-returns}
\end{figure}

\subsection{Αναγνώριση Περιορισμού}
\label{subsec:exp-arch-caveat}

Το εύρημα απαιτεί προσεκτική ερμηνεία σε σχέση με τη σύγκριση εντός της ίδιας ομάδας και τη σύγκριση μεταξύ ομάδων. Εντός της ομάδας αρχιτεκτονικής, ο παράγοντας ανταμοιβή και το σύνολο χαρακτηριστικών είναι κοινοί, οπότε η εμφανής υπεροχή της PPO-LSTM 2$\times$512 έναντι της PPO-MLP (Sharpe 0,676 έναντι 0,313) αποτελεί αξιόπιστη ένδειξη αρχιτεκτονικής υπεροχής.

Ωστόσο, κάθε διαγωνική σύγκριση μεταξύ αυτής της ομάδας και της επόμενης (σχεδιασμός ανταμοιβής) θα ήταν αναξιόπιστη, καθώς οι δύο ομάδες διαφέρουν σε πλέον του ενός παράγοντα. Αυτός ο σχεδιαστικός περιορισμός αναγνωρίζεται ρητά. Το ουσιαστικό συμπέρασμα παραμένει: η PPO-LSTM 2$\times$512 επιλέγεται ως κανονική αρχιτεκτονική για τις επόμενες ομάδες, τόσο λόγω της ανώτερης εκτός δείγματος Sharpe όσο και λόγω της θεωρητικής συνέπειας με το πλαίσιο POMDP.

% ==============================================================
\section{Σχεδιασμός Συνάρτησης Ανταμοιβής}
\label{sec:exp-reward}
% ==============================================================

\subsection{Κίνητρο}
\label{subsec:exp-reward-motivation}

Σε ένα περιβάλλον χαμηλού λόγου σήματος-θορύβου όπως τα κρυπτονομίσματα, ο πράκτορας DRL τείνει να συγκλίνει σε παθητική κατάρρευση: σταθερή παραμονή σε θέση μετρητών ή πλήρης στατική επένδυση χωρίς περιστροφή χαρτοφυλακίου. Η κεντρική υπόθεση της ομάδας σχεδιασμού ανταμοιβής είναι ότι η ανταμοιβή ενεργής απόδοσης (active return) αποτρέπει αυτή την κατάρρευση επιβάλλοντας σαφές αρνητικό σήμα σε παθητικές στρατηγικές: ο πράκτορας που απλώς ακολουθεί τον ισόποσο δείκτη αναφοράς λαμβάνει ανταμοιβή μηδέν, ακόμη και αν η αγορά ανεβαίνει.

Πέντε διαμορφώσεις συγκρίνονται με σταθερή αρχιτεκτονική PPO-LSTM 2$\times$512, σταθερό σύνολο χαρακτηριστικών και κόστος $10^{-4}$.

\subsection{Αποτελέσματα}
\label{subsec:exp-reward-results}

Ο Πίνακας~\ref{tab:g3-test} παρουσιάζει τα αποτελέσματα εκτός δείγματος για τις πέντε διαμορφώσεις ανταμοιβής.

\begin{table}[htbp]
  \centering
  \caption{Αποτελέσματα σύγκρισης συναρτήσεων ανταμοιβής (PPO-LSTM 2$\times$512, κόστος $10^{-4}$, περίοδος αξιολόγησης 2024--2025).}
  \label{tab:g3-test}
  \begin{tabular}{@{}lcccc@{}}
    \toprule
    \textbf{Συνάρτηση Ανταμοιβής} & \textbf{Sharpe} & \textbf{Μεταβλ.} & \textbf{Max DD} & \textbf{Τελικό (K)} \\
    \midrule
    PnL απόδοση                              & 0,662 & 0,699 & $-0{,}472$ & 1.234 \\
    \textbf{Ενεργή απόδοση (active return)}  & \textbf{0,778} & \textbf{0,702} & $-0{,}508$ & \textbf{1.335} \\
    Ανταμοιβή Sharpe                         & 0,035 & 0,742 & $-0{,}574$ & 783 \\
    Αυξανόμενο παράθυρο Sharpe              & 0,623 & 0,712 & $-0{,}474$ & 1.201 \\
    Ενεργή απόδοση + ποινή κύκλου εργασιών  & 0,431 & 0,764 & $-0{,}556$ & 1.036 \\
    \midrule
    Buy-and-Hold (αναφορά)                   & 0,267 & 0,754 & $-0{,}583$ & 922 \\
    \bottomrule
  \end{tabular}
\end{table}

\begin{figure}[htbp]
  \centering
  \includegraphics[width=0.85\textwidth]{../../out/final_groups/trade/g3/cumulative_returns_mean.png}
  \caption{Σωρευτικές αποδόσεις εκτός δείγματος για τις πέντε διαμορφώσεις ανταμοιβής. Η ανταμοιβή Sharpe συγκλίνει σε ζημία, ενώ η ανταμοιβή ενεργής απόδοσης υπερτερεί σαφώς.}
  \label{fig:g3-returns}
\end{figure}

\begin{figure}[htbp]
  \centering
  \includegraphics[width=0.85\textwidth]{../../out/final_groups/trade/g3/cumulative_returns_relative_bnh.png}
  \caption{Σωρευτικές αποδόσεις σχετικά με τη βασική γραμμή αγοράς-και-κράτησης. Η ανταμοιβή Sharpe εμφανίζει αρνητική σχετική απόδοση, επιβεβαιώνοντας την παθητική κατάρρευση.}
  \label{fig:g3-relative}
\end{figure}

\subsection{Ανάλυση Αποτελεσμάτων}
\label{subsec:exp-reward-analysis}

Η ανταμοιβή ενεργής απόδοσης επιτυγχάνει τη βέλτιστη Sharpe (0,778) και τελική αξία (1.335K). Τα διαγνωστικά κατανομής ενεργειών αποκαλύπτουν τον λόγο: ο πράκτορας αναπτύσσει διακριτή ανά-μέσο στρατηγική, με το XRP να διατηρεί σταθερά θετικές ενέργειες (+0,4 ως +0,6) και το TRX σταθερά αρνητικές ($-0{,}4$ ως $-0{,}6$). Αυτή η διαφοροποίηση είναι η διαγνωστική υπογραφή αποτελεσματικής ανακάλυψης σήματος και απουσιάζει από τις υπόλοιπες διαμορφώσεις.

Η ανταμοιβή Sharpe αποδίδει Sharpe 0,035, κατώτερα από τη βασική γραμμή, με τελική αξία 783K δηλαδή πραγματική ζημία. Η αποτυχία αυτή συνάδει με γνωστή αδυναμία της μεθόδου: η κλίση βαθμωτής συνάρτησης που εξαρτάται από τη στιγμιαία τυπική απόκλιση είναι αστατής όταν η πολιτική παράγει υψηλή μεταβλητότητα, ιδίως στα πρώτα στάδια εκπαίδευσης. Το αυξανόμενο παράθυρο Sharpe (0,623) μετριάζει αυτή την αστάθεια χρησιμοποιώντας μεγαλύτερο ιστορικό παράθυρο, αλλά παραμένει κατώτερο από την ενεργή απόδοση.

Η ποινή κύκλου εργασιών (0,431) εισάγει αντίδραση που αποθαρρύνει αρκετά τις συναλλαγές ώστε να ωθεί τον πράκτορα μερικώς προς παθητική συμπεριφορά, ενώ η ανταμοιβή PnL (0,662) αποδίδει ικανοποιητικά αλλά χωρίς σήμα σχετικής απόδοσης έναντι benchmark. Η ανταμοιβή ενεργής απόδοσης επιλέγεται ως κανονική διαμόρφωση για τις επόμενες ομάδες.

% ==============================================================
\section{Σύνολο Χαρακτηριστικών Εισόδου}
\label{sec:exp-features}
% ==============================================================

\subsection{Κίνητρο και Πειραματικός Σχεδιασμός}
\label{subsec:exp-features-design}

Αφού καθοριστεί η βέλτιστη αρχιτεκτονική και συνάρτηση ανταμοιβής, η ομάδα χαρακτηριστικών εξετάζει αν εμπλουτισμένα τεχνικά χαρακτηριστικά βελτιώνουν τη γενίκευση. Πέντε preset χαρακτηριστικών (A, B, C, D, E) συγκρίνονται, που αποτελούν παραλλαγές 24-ωρων τεχνικών δεικτών και δεικτών ορμής. Το preset C είναι η διαμόρφωση που χρησιμοποιήθηκε στην ομάδα σχεδιασμού ανταμοιβής, οπότε χρησιμεύει και ως εσωτερική αναφορά.

\subsection{Αποτελέσματα}
\label{subsec:exp-features-results}

Ο Πίνακας~\ref{tab:g4-test} παρουσιάζει τα αποτελέσματα εκτός δείγματος.

\begin{table}[htbp]
  \centering
  \caption{Σύγκριση preset χαρακτηριστικών εισόδου (PPO-LSTM 2$\times$512, ανταμοιβή ενεργής απόδοσης, κόστος $10^{-4}$, περίοδος αξιολόγησης 2024--2025).}
  \label{tab:g4-test}
  \begin{tabular}{@{}lcccc@{}}
    \toprule
    \textbf{Preset} & \textbf{Sharpe} & \textbf{Μεταβλ.} & \textbf{Max DD} & \textbf{Τελικό (K)} \\
    \midrule
    Preset A & 0,559 & 0,642 & $-0{,}529$ & 1.158 \\
    Preset B & 0,303 & 0,645 & $-0{,}537$ & 987 \\
    \textbf{Preset C} & \textbf{0,778} & \textbf{0,702} & $-0{,}508$ & \textbf{1.335} \\
    Preset D & 0,673 & 0,733 & $-0{,}555$ & 1.242 \\
    Preset E & 0,472 & 0,704 & $-0{,}581$ & 1.084 \\
    \midrule
    Buy-and-Hold (αναφορά) & 0,267 & 0,754 & $-0{,}583$ & 922 \\
    \bottomrule
  \end{tabular}
\end{table}

\begin{figure}[htbp]
  \centering
  \includegraphics[width=0.85\textwidth]{../../out/final_groups/trade/g4/cumulative_returns_mean.png}
  \caption{Σωρευτικές αποδόσεις εκτός δείγματος για τα πέντε preset χαρακτηριστικών. Το preset C διατηρεί σαφή πλεονέκτημα, ενώ η κατάταξη των άλλων preset δεν είναι μονοτονική.}
  \label{fig:g4-returns}
\end{figure}

\subsection{Ανάλυση}
\label{subsec:exp-features-analysis}

Το preset C επιτυγχάνει τη βέλτιστη Sharpe (0,778) και τελική αξία (1.335K), επιβεβαιώνοντας τη διαμόρφωση της ομάδας σχεδιασμού ανταμοιβής. Ωστόσο, η κατάταξη των άλλων preset δεν είναι μονοτονική ούτε προβλέψιμη από την πολυπλοκότητα των χαρακτηριστικών: το preset B, που παρέχει ένα ενδιάμεσο σύνολο χαρακτηριστικών, αποδίδει χειρότερα από το απλούστερο preset A. Αυτό υποδηλώνει ότι ο αριθμός χαρακτηριστικών δεν είναι από μόνος του ο καθοριστικός παράγοντας και ότι η σύνθεση των χαρακτηριστικών παίζει σημαντικό ρόλο. Όλα τα preset υπερτερούν της βασικής γραμμής, επιβεβαιώνοντας ότι οι τεχνικοί δείκτες παρέχουν χρήσιμο σήμα.

% ==============================================================
\section{Ευαισθησία σε Κόστος Συναλλαγής}
\label{sec:exp-fees}
% ==============================================================

\subsection{Κίνητρο και Αποτελέσματα}
\label{subsec:exp-fees-results}

Το κόστος συναλλαγής αποτελεί κρίσιμη πρακτική παράμετρο: ακόμη και μικρές διαφορές στη λεπτομέρεια του brokers (maker/taker fees, slippage) μπορούν να αλλάξουν ριζικά τη βιωσιμότητα μιας στρατηγικής υψηλού κύκλου εργασιών. Η ομάδα ευαισθησίας κόστους συγκρίνει τέσσερα επίπεδα προμήθειας με σταθερή αρχιτεκτονική PPO-LSTM και ανταμοιβή ενεργής απόδοσης. Τα αποτελέσματα παρουσιάζονται στον Πίνακα~\ref{tab:g5-test}.

\begin{table}[htbp]
  \centering
  \caption{Ευαισθησία εκτός δείγματος απόδοσης σε διάφορα επίπεδα κόστους συναλλαγής (PPO-LSTM, ανταμοιβή ενεργής απόδοσης, περίοδος αξιολόγησης 2024--2025).}
  \label{tab:g5-test}
  \begin{tabular}{@{}lcccc@{}}
    \toprule
    \textbf{Κόστος Συναλλαγής} & \textbf{Sharpe} & \textbf{Μεταβλ.} & \textbf{Max DD} & \textbf{Τελικό (K)} \\
    \midrule
    $10^{-3}$ (0,1\%)  & 0,200 & 0,795 & $-0{,}574$ & 858 \\
    \textbf{$10^{-4}$ (0,01\%)} & \textbf{0,778} & \textbf{0,702} & $-0{,}508$ & \textbf{1.335} \\
    $2\times10^{-5}$   & 0,876 & 0,642 & $-0{,}461$ & 1.411 \\
    $10^{-5}$          & 0,880 & 0,635 & $-0{,}455$ & 1.411 \\
    \midrule
    Buy-and-Hold (αναφορά) & 0,267 & 0,754 & $-0{,}583$ & 922 \\
    \bottomrule
  \end{tabular}
\end{table}

\begin{figure}[htbp]
  \centering
  \includegraphics[width=0.85\textwidth]{../../out/final_groups/trade/g5/cumulative_returns_mean.png}
  \caption{Σωρευτικές αποδόσεις εκτός δείγματος για διάφορα επίπεδα κόστους συναλλαγής. Το κόστος $10^{-3}$ μειώνει δραστικά την απόδοση, ενώ τα χαμηλότερα κόστη βελτιώνουν το αποτέλεσμα αλλά με φθίνουσα οριακή αξία.}
  \label{fig:g5-returns}
\end{figure}

\subsection{Ανάλυση}
\label{subsec:exp-fees-analysis}

Το κόστος $10^{-3}$ προκαλεί δραστική υποβάθμιση: Sharpe 0,200, οριακά άνω της βασικής γραμμής (0,267), με τελική αξία 858K κάτω του αρχικού κεφαλαίου. Η τριψήφια μείωση Sharpe (από 0,778 σε 0,200) επιβεβαιώνει ότι ο πράκτορας εκμεταλλεύεται ευκαιρίες υψηλής συχνότητας που εξαλείφονται από το υψηλότερο κόστος. Το $10^{-4}$ αποτελεί το πρακτικό σημείο ισορροπίας: αντικατοπτρίζει ρεαλιστικές συνθήκες αγοράς για ρευστά κρυπτονομίσματα με maker/taker fees. Τα χαμηλότερα κόστη ($2\times10^{-5}$, $10^{-5}$) βελτιώνουν ελαφρώς (Sharpe 0,876--0,880) αλλά υπερεκτιμούν τις πραγματικές συνθήκες αγοράς. Το $10^{-4}$ επιβεβαιώνεται ως κανονικό επίπεδο κόστους.

% ==============================================================
\section{Έλεγχος Δράσεων και Έκθεσης}
\label{sec:exp-actions}
% ==============================================================

\subsection{Κίνητρο και Αποτελέσματα}
\label{subsec:exp-actions-results}

Ο τρόπος με τον οποίο οι ενέργειες του πράκτορα αντιστοιχίζονται σε πραγματικές αλλαγές χαρτοφυλακίου επηρεάζει τόσο τη μαθησιακή δυναμική όσο και τον κίνδυνο υπερβολικής έκθεσης σε μεμονωμένα μέσα. Η ομάδα ελέγχου δράσεων συγκρίνει τέσσερις μηχανισμούς: ομαλή αντιστοίχιση με βάση μετοχές, ομαλή αντιστοίχιση με βάση χαρτοφυλάκιο, κανονική αντιστοίχιση μέγιστης έκθεσης (η ήδη χρησιμοποιούμενη), και αντιστοίχιση χωρίς υπολανθάνουσα μόχλευση. Τα αποτελέσματα παρουσιάζονται στον Πίνακα~\ref{tab:g6-test}.

\begin{table}[htbp]
  \centering
  \caption{Σύγκριση μηχανισμών αντιστοίχισης δράσεων και ελέγχου έκθεσης (PPO-LSTM, ενεργή απόδοση, κόστος $10^{-4}$, περίοδος αξιολόγησης 2024--2025).}
  \label{tab:g6-test}
  \begin{tabular}{@{}lcccc@{}}
    \toprule
    \textbf{Μηχανισμός Δράσεων} & \textbf{Sharpe} & \textbf{Μεταβλ.} & \textbf{Max DD} & \textbf{Τελικό (K)} \\
    \midrule
    Ομαλή αντιστοίχιση (μετοχές)       & 0,445 & 0,865 & $-0{,}570$ & 1.010 \\
    Ομαλή αντιστοίχιση (χαρτοφυλάκιο) & 0,278 & 0,713 & $-0{,}562$ & 946 \\
    \textbf{Μέγιστη έκθεση (κανονική)} & \textbf{0,778} & \textbf{0,702} & $-0{,}508$ & \textbf{1.335} \\
    Χωρίς υπολανθάνουσα μόχλευση      & 0,430 & 0,791 & $-0{,}591$ & 1.026 \\
    \midrule
    Buy-and-Hold (αναφορά)             & 0,267 & 0,754 & $-0{,}583$ & 922 \\
    \bottomrule
  \end{tabular}
\end{table}

\begin{figure}[htbp]
  \centering
  \includegraphics[width=0.85\textwidth]{../../out/final_groups/trade/g6/cumulative_returns_mean.png}
  \caption{Σωρευτικές αποδόσεις εκτός δείγματος για τους τέσσερις μηχανισμούς αντιστοίχισης δράσεων. Η κανονική αντιστοίχιση μέγιστης έκθεσης υπερτερεί σαφώς.}
  \label{fig:g6-returns}
\end{figure}

Η κανονική αντιστοίχιση μέγιστης έκθεσης επιτυγχάνει Sharpe 0,778, υπερτερώντας και των δύο εναλλακτικών ομαλής αντιστοίχισης. Η ομαλή αντιστοίχιση βάσει μετοχών (0,445) περιορίζει την ικανότητα του πράκτορα να πραγματοποιεί μεγάλες γρήγορες αναδιαρθρώσεις, κάτι που αποδεικνύεται επιβλαβές στο τρέχον περιβάλλον. Η αντιστοίχιση χωρίς υπολανθάνουσα μόχλευση (0,430) εισάγει πρόσθετους περιορισμούς που μειώνουν ελαφρώς την απόδοση. Η κανονική αντιστοίχιση επιβεβαιώνεται ως η αποτελεσματικότερη.

% ==============================================================
\section{Μήκος Επεισοδίου και Παράλληλα Περιβάλλοντα}
\label{sec:exp-episodes}
% ==============================================================

\subsection{Κίνητρο και Αποτελέσματα}
\label{subsec:exp-episodes-results}

Το μήκος επεισοδίου καθορίζει πόσο μακρύ χρονικό ορίζοντα έχει η LSTM για να αναπτύξει χρονική στρατηγική εντός κάθε rollout. Μεγαλύτερα επεισόδια παρέχουν πλουσιότερο χρονικό πλαίσιο αλλά μειώνουν την ποικιλομορφία των δειγμάτων. Τέσσερα συνδυαστικά ζεύγη (βήματα ανά επεισόδιο, αριθμός παράλληλων περιβαλλόντων) αξιολογούνται διατηρώντας σταθερό τον συνολικό αριθμό βημάτων ανά εποχή. Τα αποτελέσματα παρουσιάζονται στον Πίνακα~\ref{tab:g7-test}.

\begin{table}[htbp]
  \centering
  \caption{Επίδραση μήκους επεισοδίου στην εκτός δείγματος απόδοση (PPO-LSTM, ενεργή απόδοση, κόστος $10^{-4}$, περίοδος αξιολόγησης 2024--2025).}
  \label{tab:g7-test}
  \begin{tabular}{@{}llcccc@{}}
    \toprule
    \textbf{Βήματα} & \textbf{Περιβάλλοντα} & \textbf{Sharpe} & \textbf{Μεταβλ.} & \textbf{Max DD} & \textbf{Τελικό (K)} \\
    \midrule
    \textbf{2.000} & \textbf{16}   & \textbf{0,778} & \textbf{0,702} & $-0{,}508$ & \textbf{1.335} \\
    720            & 45            & 0,470          & 0,706          & $-0{,}554$ & 1.082 \\
    168            & 191           & 0,388          & 0,697          & $-0{,}568$ & 1.026 \\
    24             & 1.334         & 0,257          & 0,859          & $-0{,}628$ & 866 \\
    \midrule
    Buy-and-Hold (αναφορά) & -- & 0,267 & 0,754 & $-0{,}583$ & 922 \\
    \bottomrule
  \end{tabular}
\end{table}

\begin{figure}[htbp]
  \centering
  \includegraphics[width=0.85\textwidth]{../../out/final_groups/trade/g7/cumulative_returns_mean.png}
  \caption{Σωρευτικές αποδόσεις εκτός δείγματος για τα τέσσερα ζεύγη μήκους επεισοδίου και αριθμού παράλληλων περιβαλλόντων. Η μονοτονική βελτίωση με το μήκος επεισοδίου επιβεβαιώνει τη σημασία του χρονικού ορίζοντα.}
  \label{fig:g7-returns}
\end{figure}

Η κατάταξη είναι μονοτονική: η Sharpe αυξάνεται με το μήκος επεισοδίου από 0,257 (24 βήματα) έως 0,778 (2.000 βήματα). Τα πολύ σύντομα επεισόδια (24 βήματα, 1.334 περιβάλλοντα) αποδίδουν χειρότερα ακόμη και από τη βασική γραμμή: η LSTM δεν έχει αρκετό χρόνο εντός του επεισοδίου για να αναπτύξει συνεκτική χρονική στρατηγική. Η μονοτονική βελτίωση επιβεβαιώνει ότι ο χρονικός ορίζοντας εντός του rollout είναι κρίσιμος για τη συγκεκριμένη αρχιτεκτονία LSTM. Το ζεύγος 2.000 βημάτων με 16 περιβάλλοντα επιλέγεται ως κανονική παραμετροποίηση.

% ==============================================================
\section{Περιορισμοί Ρευστότητας μέσω Σκληρής Ποινής}
\label{sec:exp-hard-liq}
% ==============================================================

\subsection{Κίνητρο και Πειραματικός Σχεδιασμός}
\label{subsec:exp-hard-liq-design}

Αφού εδραιωθεί η κανονική διαμόρφωση (PPO-LSTM 2$\times$512, ενεργή απόδοση, κόστος $10^{-4}$, 2.000 βήματα), ο δεύτερος ερευνητικός άξονας εξετάζει αν η επιβολή περιορισμών ρευστότητας κατά την εκπαίδευση βελτιώνει το προφίλ κινδύνου-απόδοσης. Η πρώτη ομάδα ρευστότητας δοκιμάζει τρεις μορφές σκληρής ποινής: ένα soft-start buffer (στόχος 10\% σε μετρητά με ποινή που ενεργοποιείται σταδιακά), ένα σκληρό buffer (στόχος 20\% με σταθερή ποινή), και κανονιστικοποίηση L2 βαρών ως εναλλακτική μορφή ρύθμισης. Τα αποτελέσματα παρουσιάζονται στον Πίνακα~\ref{tab:g8_1-test}.

\begin{table}[htbp]
  \centering
  \caption{Αποτελέσματα σκληρών περιορισμών ρευστότητας (PPO-LSTM, ενεργή απόδοση, κόστος $10^{-4}$, περίοδος αξιολόγησης 2024--2025). Η ανεξάρτητη αναφορά ενεργής απόδοσης χωρίς περιορισμό συμπεριλαμβάνεται για άμεση σύγκριση.}
  \label{tab:g8_1-test}
  \begin{tabular}{@{}lcccc@{}}
    \toprule
    \textbf{Διαμόρφωση} & \textbf{Sharpe} & \textbf{Μεταβλ.} & \textbf{Max DD} & \textbf{Τελικό (K)} \\
    \midrule
    Ενεργή απόδοση (αναφορά)    & 0,778 & 0,702 & $-0{,}508$ & 1.335 \\
    \midrule
    Soft-start buffer 10\%      & 0,693 & 0,492 & $-0{,}325$ & 1.236 \\
    Σκληρό buffer 20\%          & 0,630 & 0,287 & $-0{,}251$ & 1.145 \\
    Κανονιστικοποίηση L2 βαρών  & 0,271 & 0,752 & $-0{,}576$ & 925 \\
    \midrule
    Buy-and-Hold (αναφορά)      & 0,267 & 0,754 & $-0{,}583$ & 922 \\
    \bottomrule
  \end{tabular}
\end{table}

\begin{figure}[htbp]
  \centering
  \includegraphics[width=0.85\textwidth]{../../out/final_groups/trade/g8_1/cumulative_returns_mean.png}
  \caption{Σωρευτικές αποδόσεις εκτός δείγματος για τους σκληρούς περιορισμούς ρευστότητας. Ο soft-start buffer μειώνει τη μεταβλητότητα ενώ διατηρεί ικανοποιητική απόδοση. Η κανονιστικοποίηση L2 δεν υπερτερεί της βασικής γραμμής.}
  \label{fig:g8_1-returns}
\end{figure}

\subsection{Ανάλυση}
\label{subsec:exp-hard-liq-analysis}

Τα αποτελέσματα αποκαλύπτουν ένα σαφές ανταλλακτικό: ο soft-start buffer μειώνει σημαντικά τη μεταβλητότητα (από 0,702 σε 0,492) και τη μέγιστη πτώση (από $-50{,}8\%$ σε $-32{,}5\%$) με μέτριο μόνο κόστος σε Sharpe (0,693 έναντι 0,778) και τελική αξία (1.236K έναντι 1.335K). Το σκληρό buffer 20\% μειώνει περαιτέρω τη μεταβλητότητα (0,287) αλλά με επίσης χαμηλότερη τελική αξία (1.145K). Η κανονιστικοποίηση L2 αποτυγχάνει να προσφέρει πλεονέκτημα έναντι της βασικής γραμμής (Sharpe 0,271, 925K), επιβεβαιώνοντας ότι η ρύθμιση των παραμέτρων του νευρωνικού δικτύου δεν αντικαθιστά τους εφαρμοσμένους περιορισμούς ρευστότητας στο περιβάλλον.

Το σημαντικό εύρημα είναι ότι ο soft-start buffer μειώνει τη μεταβλητότητα χωρίς να εισάγει νέες μορφές κατάρρευσης. Αυτή η ιδιότητα θα εξεταστεί περαιτέρω στην ομάδα CPCV βελτιστοποίησης (Ενότητα~\ref{sec:exp-cpcv}).

% ==============================================================
\section{Λαγκρανζιανή Διαμόρφωση Ρευστότητας}
\label{sec:exp-lagrangian}
% ==============================================================

\subsection{Πειραματικός Σχεδιασμός}
\label{subsec:exp-lagrangian-design}

Ο δεύτερος άξονας ρευστότητας επεκτείνεται σε Λαγκρανζιανή διατύπωση CMDP \cite{altman1999cmdp}, όπου ο περιορισμός ρευστότητας εισάγεται μέσω εκμαθημένου πολλαπλασιαστή $\lambda$ αντί σκληρής ποινής. Πέντε διαμορφώσεις Λαγκρανζιανής χαλάρωσης συγκρίνονται: ομαλή κυρτή έλλειψη ρευστότητας, EMA αναλογία μετρητών, Λαγκρανζιανό buffer 20\%, Lagrangian liquify (10\% ταμείο όταν θέσεις $\geq 90\%$), και διπλόπλευρη ζώνη μετρητών 15/80\%. Τα αποτελέσματα παρουσιάζονται στον Πίνακα~\ref{tab:g8_2-test}.

\begin{table}[htbp]
  \centering
  \caption{Αποτελέσματα Λαγκρανζιανών διαμορφώσεων ρευστότητας (PPO-LSTM, ενεργή απόδοση, κόστος $10^{-4}$, περίοδος αξιολόγησης 2024--2025). Σημείωση: οι τιμές Sharpe$>$1,0 με $\sigma < 0{,}15$ υποδεικνύουν πιθανή γωνιακή λύση (βλ.\ Ενότητα~\ref{subsec:exp-ema-paradox}).}
  \label{tab:g8_2-test}
  \begin{tabular}{@{}lcccc@{}}
    \toprule
    \textbf{Διαμόρφωση} & \textbf{Sharpe} & \textbf{Μεταβλ.} & \textbf{Max DD} & \textbf{Τελικό (K)} \\
    \midrule
    Ενεργή απόδοση (αναφορά)         & 0,778 & 0,702 & $-0{,}508$ & 1.335 \\
    \midrule
    Λαγκρ.\ ομαλή κυρτή έλλειψη      & 1,026 & 0,131 & $-0{,}117$ & 1.129 \\
    Λαγκρ.\ EMA αναλογία μετρητών$^*$ & 1,184 & 0,084 & $-0{,}042$ & 1.097 \\
    Λαγκρ.\ buffer 20\%               & 0,754 & 0,342 & $-0{,}243$ & 1.212 \\
    Λαγκρ.\ liquify 10\% στο 90\%    & 0,722 & 0,447 & $-0{,}313$ & 1.240 \\
    Διπλόπλευρη ζώνη 15/80\%         & 0,372 & 0,396 & $-0{,}302$ & 1.069 \\
    \midrule
    Buy-and-Hold (αναφορά)           & 0,267 & 0,754 & $-0{,}583$ & 922 \\
    \bottomrule
    \multicolumn{5}{@{}l@{}}{\small $^*$Γωνιακή λύση -- βλ.\ Ενότητα~\ref{subsec:exp-ema-paradox}}
  \end{tabular}
\end{table}

\begin{figure}[htbp]
  \centering
  \includegraphics[width=0.85\textwidth]{../../out/final_groups/trade/g8_2/cumulative_returns_mean.png}
  \caption{Σωρευτικές αποδόσεις εκτός δείγματος για τις Λαγκρανζιανές διαμορφώσεις ρευστότητας. Η υψηλή Sharpe ratio του πράκτορα EMA αναλογίας μετρητών δεν αντικατοπτρίζει γνήσια επενδυτική συμπεριφορά (βλ.\ κείμενο).}
  \label{fig:g8_2-returns}
\end{figure}

\subsection{Το Παράδοξο της EMA: Υψηλός Sharpe ως Αποτέλεσμα Κατάρρευσης}
\label{subsec:exp-ema-paradox}

Ο πράκτορας EMA αναλογίας μετρητών εμφανίζει στον Πίνακα~\ref{tab:g8_2-test} τη μέγιστη Sharpe ratio (1,184) και ελάχιστη μέγιστη πτώση ($-4{,}2\%$). Μια επιφανειακή ανάγνωση των μετρικών θα οδηγούσε στο συμπέρασμα ότι αυτός είναι ο κατ' εξοχήν επιτυχημένος πράκτορας της ομάδας. Η διαγνωστική ανάλυση συμπεριφοράς αποκαλύπτει ότι αυτό το συμπέρασμα είναι εσφαλμένο.

\textbf{Διαγνωστική συμπεριφορά.} Κατά τη διάρκεια της περιόδου αξιολόγησης, τα μετρητά παραμένουν μεταξύ 750K και 1,1M χωρίς σημαντική διακύμανση. Το ποσοστό κατεχόμενων θέσεων σε μέσα δεν υπερβαίνει το 15\%. Τα μέσα κατανομής ενεργειών για όλα τα δέκα περιουσιακά στοιχεία είναι αρνητικά, συγκεντρωμένα στο διάστημα $[-1{,}0, -3{,}0]$, με το ETH να διατηρεί τιμή κοντά στο $-3{,}0$ σταθερά. Αυτό είναι η διαγνωστική υπογραφή μιας πολιτικής «πώλησε πάντα, παράμεινε σε μετρητά».

\textbf{Μηχανισμός γωνιακής λύσης.} Ο EMA περιορισμός ορίζει ότι η εκθετικά κινητή μέση του ποσοστού μετρητών πρέπει να παραμένει εντός ενός στόχου. Η κράτηση μετρητών ικανοποιεί τετριμμένα αυτό τον περιορισμό: αν ο πράκτορας απλώς αποφεύγει κάθε αγορά, η EMA αναλογία μετρητών διατηρείται σταθερά στο μέγιστο επιτρεπτό επίπεδο. Η Λαγκρανζιανή δυναμική ενισχύει αυτό το κίνητρο: ο πολλαπλασιαστής $\lambda$ αυξάνεται όταν ο περιορισμός παραβιάζεται, αλλά η κράτηση μετρητών δεν τον παραβιάζει ποτέ, οπότε $\lambda$ δεν αποθαρρύνει ποτέ αυτή τη συμπεριφορά. Ο πράκτορας έχει ανακαλύψει μια τετριμμένη γωνιακή λύση που είναι συνεκτική αλλά επενδυτικά αδρανής.

\textbf{Γιατί η Sharpe ratio παραπλανά.} Η μεταβλητότητα του πράκτορα EMA (0,084) είναι μόλις το 12\% της μεταβλητότητας του πράκτορα ενεργής απόδοσης (0,702). Με τόσο χαμηλή μεταβλητότητα, ακόμη και μια μέτρια θετική drift αποφέρει υψηλό Sharpe. Η τελική αξία (1.097K) είναι ωστόσο οριακά άνω της βασικής γραμμής (922K) και κατώτερη από τον πράκτορα ενεργής απόδοσης (1.335K). Ο Sharpe 1,184 δεν αντικατοπτρίζει ανώτερη ανακάλυψη αρβιτράζ αλλά αδρανή συσσώρευση μετρητών υπό ευνοϊκές συνθήκες αγοράς.

Η πλήρης ομαλή κυρτή έλλειψη (Sharpe 1,026, Μεταβλ.\ 0,131) παρουσιάζει παρόμοιο, αν και λιγότερο έντονο, μοτίβο: εξαιρετικά χαμηλή μεταβλητότητα σε συνδυασμό με τελική αξία (1.129K) μικρότερη από τον πράκτορα ενεργής απόδοσης χωρίς περιορισμό. Αυτές οι γωνιακές λύσεις αναδεικνύουν ότι η διαγνωστική ανάλυση συμπεριφοράς είναι αναγκαία συμπλήρωση της αξιολόγησης μετρικών και δεν αντικαθίσταται από αυτήν.

Οι πράκτορες Λαγκρανζιανού buffer 20\% (0,754, 1.212K) και Lagrangian liquify (0,722, 1.240K) αποδίδουν με γνήσια επενδυτική συμπεριφορά, υπερτερώντας της βασικής γραμμής αλλά υπολειπόμενοι της ανεξάρτητης αναφοράς ενεργής απόδοσης (0,778, 1.335K). Η διπλόπλευρη ζώνη 15/80\% (0,372) αποδίδει χαμηλά: ο ταυτόχρονος περιορισμός ανώτατου και κατώτατου ορίου μετρητών δημιουργεί αντικρουόμενα κίνητρα που ο πράκτορας δεν μπορεί να βελτιστοποιήσει αποτελεσματικά χωρίς CPCV βελτιστοποίηση.

% ==============================================================
\section{Βελτιστοποίηση Υπερπαραμέτρων μέσω CPCV}
\label{sec:exp-cpcv}
% ==============================================================

\subsection{Κίνητρο και Πειραματικός Σχεδιασμός}
\label{subsec:exp-cpcv-design}

Τα αποτελέσματα των προηγούμενων ομάδων χρησιμοποίησαν σταθερές υπερπαραμέτρους. Η ομάδα CPCV βελτιστοποίησης εξετάζει κατά πόσο η συστηματική αναζήτηση υπερπαραμέτρων μέσω Combinatorial Purged Cross-Validation \cite{lopezdeprado2018} και Optuna TPE \cite{akiba2019optuna} αναδεικνύει καλύτερες παραμετροποιήσεις χωρίς υπερπροσαρμογή στην περίοδο εκπαίδευσης. Πέντε διαμορφώσεις βελτιστοποιούνται: ενεργή απόδοση (εσωτερική αναφορά), soft-start buffer 10\%, Λαγκρανζιανό buffer 20\%, Lagrangian liquify 10\% στο 90\% θέσεων, και διπλόπλευρη ζώνη 5/90\%. Τα αποτελέσματα παρουσιάζονται στον Πίνακα~\ref{tab:g9-test}.

\begin{table}[htbp]
  \centering
  \caption{Αποτελέσματα CPCV/Optuna βελτιστοποίησης ανά διαμόρφωση (PPO-LSTM, κόστος $10^{-4}$, \texttt{equal\_cash\_holds}, περίοδος αξιολόγησης 2024--2025).}
  \label{tab:g9-test}
  \begin{tabular}{@{}lcccc@{}}
    \toprule
    \textbf{Διαμόρφωση} & \textbf{Sharpe} & \textbf{Μεταβλ.} & \textbf{Max DD} & \textbf{Τελικό (K)} \\
    \midrule
    Ενεργή απόδοση (αναφορά CPCV) & 0,773 & 0,702 & $-0{,}509$ & 1.330 \\
    \textbf{Soft-start buffer 10\%} & \textbf{1,026} & \textbf{0,569} & $-0{,}367$ & \textbf{1.503} \\
    Λαγκρ.\ buffer 20\%            & 0,530 & 0,573 & $-0{,}388$ & 1.143 \\
    Λαγκρ.\ liquify 10\% στο 90\%  & 0,721 & 0,450 & $-0{,}315$ & 1.240 \\
    Διπλόπλευρη ζώνη 5/90\%        & $-0{,}154$ & 0,373 & $-0{,}366$ & 885 \\
    \midrule
    Buy-and-Hold (αναφορά)         & 0,267 & 0,754 & $-0{,}583$ & 922 \\
    \bottomrule
  \end{tabular}
\end{table}

\begin{figure}[htbp]
  \centering
  \includegraphics[width=0.85\textwidth]{../../out/final_groups/trade/g9/cumulative_returns_mean.png}
  \caption{Σωρευτικές αποδόσεις εκτός δείγματος για τις CPCV-βελτιστοποιημένες διαμορφώσεις. Ο soft-start buffer αναδεικνύεται ως η καλύτερη διαμόρφωση, ενώ η διπλόπλευρη ζώνη 5/90\% καταρρέει.}
  \label{fig:g9-returns}
\end{figure}

\begin{figure}[htbp]
  \centering
  \includegraphics[width=0.85\textwidth]{../../out/final_groups/trade/g9/cumulative_returns_relative_bnh.png}
  \caption{Σωρευτικές αποδόσεις σχετικά με τη βασική γραμμή αγοράς-και-κράτησης για τις CPCV-βελτιστοποιημένες διαμορφώσεις. Ο soft-start buffer διατηρεί σταθερά θετική σχετική απόδοση.}
  \label{fig:g9-relative}
\end{figure}

\subsection{Ανάλυση}
\label{subsec:exp-cpcv-analysis}

Το πιο εντυπωσιακό εύρημα είναι η άνοδος του soft-start buffer: από Sharpe 0,693 χωρίς βελτιστοποίηση (Ενότητα~\ref{sec:exp-hard-liq}) σε Sharpe 1,026 με CPCV/Optuna βελτιστοποίηση. Ταυτόχρονα, η τελική αξία αυξάνεται από 1.236K σε 1.503K, δηλαδή η βελτιστοποίηση βελτιώνει όχι μόνο τον λόγο Sharpe αλλά και την απόλυτη απόδοση. Αυτός ο συνδυασμός (ανώτερη Sharpe ΚΑΙ ανώτερη τελική αξία από τον πράκτορα ενεργής απόδοσης χωρίς περιορισμό) καθιστά τον βελτιστοποιημένο soft-start buffer την καλύτερη διαμόρφωση καθαρής RL (χωρίς σήμα συναισθήματος) μεταξύ όλων των ομάδων.

Ο μηχανισμός είναι ο εξής: ο soft-start περιορισμός εισάγεται με χαμηλό αρχικό $\lambda$, επιτρέποντας στον πράκτορα να αναπτύξει πρώτα μια ενεργή επενδυτική πολιτική και να αποφύγει τη γωνιακή λύση. Καθώς η εκπαίδευση εξελίσσεται, ο $\lambda$ αυξάνεται, ενισχύοντας τον περιορισμό ρευστότητας σε ήδη εδραιωμένη ενεργή πολιτική. Η CPCV βελτιστοποίηση βρίσκει τη βέλτιστη παραμετροποίηση αυτής της δυναμικής ανά CPCV fold.

Η διπλόπλευρη ζώνη 5/90\% καταρρέει πλήρως (Sharpe $-0{,}154$, 885K): η ταυτόχρονη επιβολή ανώτατου ορίου (5\%) και κατώτατου ορίου (90\% θέσεων) δημιουργεί ένα κενό πολιτικής που ο πράκτορας δεν μπορεί να ικανοποιήσει. Ο Λαγκρανζιανός buffer 20\% αποδίδει μετρίως (0,530), χαμηλότερα από την CPCV αναφορά ενεργής απόδοσης (0,773). Αυτό υποδηλώνει ότι ο τύπος του περιορισμού, και όχι μόνο η βελτιστοποίηση των υπερπαραμέτρων, είναι καθοριστικός παράγοντας.

% ==============================================================
\section{Ενσωμάτωση Σήματος Συναισθήματος FinGPT}
\label{sec:exp-sentiment}
% ==============================================================

\subsection{Περιγραφή Χαρακτηριστικών}
\label{subsec:exp-sentiment-features}

Η τρίτη ερευνητική συνεισφορά εξετάζει αν η ενσωμάτωση σήματος συναισθήματος από το FinGPT \cite{yang2023fingpt} παρέχει πρόσθετη αξία, ιδίως σε συνδυασμό με Λαγκρανζιανούς περιορισμούς ρευστότητας. Τα χαρακτηριστικά εισόδου επεκτείνονται από τον κανονικό κόμβο (19 χαρακτηριστικά βάσης και ορμής) στον κόμβο \texttt{base\_momentum\_fingpt\_onehot}: για κάθε ένα από τα 10 κρυπτονομίσματα, μια ανά-μέσο κωδικοποίηση one-hot 5 κλάσεων (\texttt{strongly\_negative}, \texttt{negative}, \texttt{neutral}, \texttt{positive}, \texttt{strongly\_positive}) από αξιολόγηση FinGPT 6,5 εκατομμυρίων χρηματοοικονομικών ειδήσεων, καθώς και ένα συνολικό αγοραίο σκορ. Ο συνολικός αριθμός χαρακτηριστικών είναι 19 + 50 (one-hot) + 1 (αγοραίο σκορ) = 70 ανά βήμα.

Η ομάδα FinGPT επαναλαμβάνει τις πέντε διαμορφώσεις της ομάδας CPCV βελτιστοποίησης με τα εμπλουτισμένα χαρακτηριστικά, ενώ ο συντελεστής εντροπίας \texttt{ent\_coef} βελτιστοποιείται ελεύθερα από το Optuna αντί να είναι σταθερός.

\subsection{Αποτελέσματα}
\label{subsec:exp-sentiment-results}

Ο Πίνακας~\ref{tab:g12-test} παρουσιάζει τα αποτελέσματα εκτός δείγματος για τις πέντε διαμορφώσεις FinGPT.

\begin{table}[htbp]
  \centering
  \caption{Αποτελέσματα ομάδας FinGPT (χαρακτηριστικά \texttt{base\_momentum\_fingpt\_onehot}, CPCV/Optuna, κόστος $10^{-4}$, περίοδος αξιολόγησης 2024--2025). Ο συντελεστής \texttt{ent\_coef} βελτιστοποιήθηκε από το Optuna για κάθε διαμόρφωση.}
  \label{tab:g12-test}
  \begin{tabular}{@{}lccccc@{}}
    \toprule
    \textbf{Διαμόρφωση} & \textbf{ent\_coef} & \textbf{Sharpe} & \textbf{Μεταβλ.} & \textbf{Max DD} & \textbf{Τελικό (K)} \\
    \midrule
    Ενεργή απόδοση            & 0,098 & 0,803 & 0,499 & $-0{,}413$ & 1.306 \\
    Soft-start buffer 10\%    & 0,069 & 1,021 & 0,473 & $-0{,}364$ & 1.432 \\
    Λαγκρ.\ buffer 5\%        & 0,051 & 0,766 & 0,608 & $-0{,}461$ & 1.312 \\
    \textbf{Lagr.\ liquify 15\% στο 70\%} & \textbf{0,091} & \textbf{1,480} & \textbf{0,315} & $-0{,}232$ & \textbf{1.496} \\
    Διπλόπλευρη ζώνη 10/90\% & 0,041 & 0,468 & 0,403 & $-0{,}330$ & 1.109 \\
    \midrule
    CPCV αναφορά ενεργής απόδ. & -- & 0,773 & 0,702 & $-0{,}509$ & 1.330 \\
    Buy-and-Hold (αναφορά)    & -- & 0,267 & 0,754 & $-0{,}583$ & 922 \\
    \bottomrule
  \end{tabular}
\end{table}

\begin{figure}[htbp]
  \centering
  \includegraphics[width=0.85\textwidth]{../../out/final_groups/trade/g12/cumulative_returns_mean.png}
  \caption{Σωρευτικές αποδόσεις εκτός δείγματος για τις FinGPT-εμπλουτισμένες διαμορφώσεις. Ο πράκτορας Lagr.\ liquify 15/70\% επιτυγχάνει δραματικά ανώτερη απόδοση από όλες τις άλλες διαμορφώσεις.}
  \label{fig:g12-returns}
\end{figure}

\begin{figure}[htbp]
  \centering
  \includegraphics[width=0.85\textwidth]{../../out/final_groups/trade/g12/cumulative_returns_relative_bnh.png}
  \caption{Σωρευτικές αποδόσεις σχετικά με τη βασική γραμμή αγοράς-και-κράτησης για τις FinGPT-εμπλουτισμένες διαμορφώσεις. Ο Lagr.\ liquify 15/70\% διατηρεί αξιοσημείωτη θετική σχετική απόδοση καθ' όλη την περίοδο.}
  \label{fig:g12-relative}
\end{figure}

\subsection{Ανάλυση Αποτελεσμάτων}
\label{subsec:exp-sentiment-analysis}

Δύο παρατηρήσεις ξεχωρίζουν. Πρώτη παρατήρηση: η απλή ενσωμάτωση συναισθήματος σε πράκτορα ενεργής απόδοσης (ent\_coef=0,098, Sharpe 0,803) επιφέρει οριακή μόνο βελτίωση (+0,030 Sharpe) σε σχέση με τον αντίστοιχο CPCV-βελτιστοποιημένο πράκτορα χωρίς συναίσθημα (0,773). Αυτό αποτελεί σημαντικό αρνητικό εύρημα: το σήμα συναισθήματος, χρησιμοποιούμενο απομονωμένα σε πράκτορα ελεύθερης επένδυσης, δεν αποδίδει σημαντικό πρόσθετο alpha.

Δεύτερη παρατήρηση: ο πράκτορας Lagr.\ liquify 15\% στο 70\% θέσεων με συναίσθημα επιτυγχάνει Sharpe 1,480, την υψηλότερη τιμή μεταξύ όλων των πειραματικών ομάδων. Η βελτίωση σε σχέση με τον ίδιο τύπο περιορισμού χωρίς συναίσθημα (CPCV liquify, Sharpe 0,721) είναι +0,759 Sharpe, δηλαδή δραματική και πολύ μεγαλύτερη από τη βελτίωση που παρατηρείται στον πράκτορα ενεργής απόδοσης (+0,030). Αυτή η ασυμμετρία υποδηλώνει ότι η αξία του σήματος συναισθήματος εκδηλώνεται κυρίως μέσω αλληλεπίδρασης με τον Λαγκρανζιανό περιορισμό.

\subsection{Συμπεριφορική Ανάλυση Βέλτιστου Πράκτορα}
\label{subsec:exp-sentiment-behavioral}

Τα διαγνωστικά γραφήματα του πράκτορα Lagr.\ liquify 15/70\% με συναίσθημα αποκαλύπτουν τρία χαρακτηριστικά στοιχεία που δεν παρατηρούνται σε κανέναν άλλο πράκτορα.

\textbf{Σταθερή κατανομή BTC.} Το Bitcoin διατηρεί ποσοστό θέσης περίπου 12\% καθ' όλη την περίοδο αξιολόγησης, χωρίς έντονες διακυμάνσεις. Αυτό υποδηλώνει ότι ο πράκτορας αξιολογεί το BTC ως αποθεματικό αξίας με χαμηλή εξάρτηση από βραχυπρόθεσμο αγοραίο συναίσθημα.

\textbf{Ειδησεογραφική αντίδραση ATOM.} Το ATOM εμφανίζει επαναλαμβανόμενες αιχμές άνω του 30\% θέσης σε σύντομα χρονικά διαστήματα, ακολουθούμενες από επαναφορά. Αυτή η συμπεριφορά συνάδει με ανταπόκριση σε θετικά ειδησεογραφικά γεγονότα που καταγράφει το FinGPT σήμα.

\textbf{Ενεργή διαχείριση μετρητών.} Τα μετρητά κυμαίνονται μεταξύ 200K και 1,2M με υψηλό πλάτος ταλάντευσης, αντικατοπτρίζοντας δυναμική διαχείριση και αποφεύγοντας πλήρως τη γωνιακή λύση αδρανούς κράτησης μετρητών.

Ο βελτιστοποιημένος ent\_coef (0,091) είναι ο υψηλότερος μεταξύ των πρακτόρων της ομάδας FinGPT. Η υψηλή εντροπία διατηρεί την εξερεύνηση της πολιτικής, αναγκαία για την αξιοποίηση ενός σήματος με χρονικά μεταβλητή πληροφοριακή αξία: ένα χαμηλό ent\_coef θα οδηγούσε σε πρόωρη εκμετάλλευση αγνοώντας νέα σήματα.

\subsection{Μηχανισμός Συνέργειας}
\label{subsec:exp-sentiment-mechanism}

Η αλληλεπίδραση μεταξύ σήματος συναισθήματος και Λαγκρανζιανού περιορισμού liquify μπορεί να ερμηνευθεί ως εξής. Ο περιορισμός liquify 15\%/70\% επιβάλλει ποινή όταν τα μετρητά υπερβαίνουν το 15\% ΚΑΙ οι θέσεις είναι κάτω από το 70\%, δηλαδή αποθαρρύνει ταυτόχρονα την υπερβολική ρευστοποίηση και την επενδυτική αδράνεια. Αυτός ο περιορισμός δημιουργεί πίεση για επένδυση, αλλά δεν παρέχει πληροφορία κατεύθυνσης: ο πράκτορας πρέπει να αποφασίσει σε ποια μέσα να επενδύσει.

Το σήμα συναισθήματος παρέχει ακριβώς αυτή την πληροφορία, δηλαδή ποια μέσα έχουν θετική ειδησεογραφική δυναμική. Η συνδυαστική επίδραση είναι ότι ο πράκτορας βρίσκει το βέλτιστο ανταλλακτικό: επενδύει (για να αποφύγει την ποινή liquify) στα μέσα με θετικό συναίσθημα (για να μεγιστοποιήσει την ανταμοιβή). Χωρίς το σήμα συναισθήματος, ο ίδιος περιορισμός αναγκάζει επένδυση χωρίς κατεύθυνση, με χαμηλότερο αποτέλεσμα. Χωρίς τον περιορισμό, το σήμα δεν επαρκεί να διαφοροποιήσει δραστικά τη συμπεριφορά.

% ==============================================================
\section{Συνολική Σύγκριση και Συμπεράσματα}
\label{sec:exp-synthesis}
% ==============================================================

\subsection{Συγκεντρωτικός Πίνακας Αποτελεσμάτων}
\label{subsec:exp-final-table}

Ο Πίνακας~\ref{tab:final-synthesis} συγκεντρώνει τον βέλτιστο πράκτορα από κάθε πειραματική ομάδα, αναδεικνύοντας την εξέλιξη των μετρικών διαμέσου των ένδεκα ομάδων. Για ορισμένες ομάδες (χαρακτηριστικά, κόστος, δράσεις, μήκος επεισοδίου), η καλύτερη διαμόρφωση ταυτίζεται με την κανονική ρύθμιση, επιβεβαιώνοντας τις σχεδιαστικές επιλογές.

\begin{table}[htbp]
  \centering
  \caption{Συγκεντρωτικά αποτελέσματα εκτός δείγματος: βέλτιστος πράκτορας ανά πειραματική ομάδα (περίοδος αξιολόγησης 22/11/2024--09/11/2025). Σύμβολα: Vol = ετησιοποιημένη μεταβλητότητα, DD = μέγιστη πτώση.}
  \label{tab:final-synthesis}
  \small
  \begin{tabular}{@{}llcccc@{}}
    \toprule
    \textbf{Ομάδα} & \textbf{Βέλτιστη Διαμόρφωση} & \textbf{Sharpe} & \textbf{Vol} & \textbf{Max DD} & \textbf{Τελικό (K)} \\
    \midrule
    Αναφορά         & Buy-and-Hold            & 0,267 & 0,754 & $-0{,}583$ & 922 \\
    \midrule
    Βιβλιοθήκες    & ElegantRL PPO-MLP       & 0,349 & 0,752 & $-0{,}561$ & 980 \\
    Αρχιτεκτονική  & PPO-LSTM 2$\times$512 PnL & 0,676 & 0,687 & $-0{,}467$ & 1.246 \\
    Ανταμοιβή      & Active return           & 0,778 & 0,702 & $-0{,}508$ & 1.335 \\
    Χαρακτηριστικά & Preset C (= κανονικό)   & 0,778 & 0,702 & $-0{,}508$ & 1.335 \\
    Κόστος         & $10^{-4}$ (κανονικό)   & 0,778 & 0,702 & $-0{,}508$ & 1.335 \\
    Δράσεις        & Μέγιστη έκθεση (κανονικό) & 0,778 & 0,702 & $-0{,}508$ & 1.335 \\
    Μήκος επεισ.   & 2.000 βήματα (κανονικό) & 0,778 & 0,702 & $-0{,}508$ & 1.335 \\
    Σκληρή ρευστ.  & Soft-start buffer 10\%  & 0,693 & 0,492 & $-0{,}325$ & 1.236 \\
    Λαγκρ.\ ρευστ.  & Lagr.\ buffer 20\%$^\dagger$ & 0,754 & 0,342 & $-0{,}243$ & 1.212 \\
    CPCV βελτ.     & Soft-start (CPCV-βελτ.) & 1,026 & 0,569 & $-0{,}367$ & 1.503 \\
    FinGPT         & Lagr.\ liquify 15/70\%  & \textbf{1,480} & \textbf{0,315} & $-0{,}232$ & \textbf{1.496} \\
    \bottomrule
    \multicolumn{6}{@{}l@{}}{\small $^\dagger$Εξαιρείται ο EMA (Sharpe 1,184) λόγω γωνιακής λύσης (βλ.\ Ενότητα~\ref{subsec:exp-ema-paradox})}
  \end{tabular}
\end{table}

\subsection{Πρόοδος μεταξύ Ομάδων}
\label{subsec:exp-progression}

Ο Πίνακας~\ref{tab:final-synthesis} αποτυπώνει μια διαδρομή τριών σταδίων. Στο πρώτο στάδιο (βιβλιοθήκες, αρχιτεκτονική, ανταμοιβή), η Sharpe αυξάνεται από 0,267 (βασική γραμμή) σε 0,778, κυρίως μέσω επιλογής αρχιτεκτονικής (PPO-LSTM) και συνάρτησης ανταμοιβής (ενεργή απόδοση). Οι ομάδες χαρακτηριστικών, κόστους, δράσεων και μήκους επεισοδίου επιβεβαιώνουν ότι οι κανονικές επιλογές είναι βέλτιστες και δεν απαιτούν τροποποίηση.

Στο δεύτερο στάδιο (περιορισμοί ρευστότητας), η προσθήκη soft-start buffer μειώνει τη μεταβλητότητα σε κόστος μικρής μείωσης Sharpe (0,693), αλλά η CPCV βελτιστοποίηση αντιστρέφει αυτή την τάση και επιτυγχάνει Sharpe 1,026 με ταυτόχρονα υψηλότερη τελική αξία (1.503K). Αυτό δείχνει ότι οι Λαγκρανζιανοί περιορισμοί δεν εισάγουν απαραίτητα ανταλλακτικό μεταξύ Sharpe και απόδοσης, αρκεί η παραμετροποίηση να βελτιστοποιηθεί κατάλληλα.

Στο τρίτο στάδιο (FinGPT), η προσθήκη σήματος συναισθήματος σε συνδυασμό με Λαγκρανζιανό περιορισμό liquify οδηγεί σε Sharpe 1,480, το υψηλότερο σε όλες τις ομάδες. Αξιοσημείωτα, η τελική αξία (1.496K) είναι ελαφρώς χαμηλότερη από τον βέλτιστο CPCV soft-start (1.503K), αλλά η υψηλότερη Sharpe συνοδεύεται από χαμηλότερη μεταβλητότητα (0,315 έναντι 0,569), δηλαδή η βελτίωση εντοπίζεται στο προφίλ κινδύνου και όχι μόνο στην απόλυτη απόδοση.

Σημαντικό εύρημα για τη σύγκριση είναι ότι η Sharpe 1,184 του πράκτορα EMA αναλογίας μετρητών εξαιρείται από τον Πίνακα~\ref{tab:final-synthesis} ως βέλτιστος ομάδας, γιατί αντικατοπτρίζει γωνιακή λύση αδρανούς κράτησης μετρητών (Ενότητα~\ref{subsec:exp-ema-paradox}) και όχι γνήσια ενεργή διαχείριση.

\subsection{Συνολικά Συμπεράσματα}
\label{subsec:exp-final-conclusions}

Τα πειραματικά αποτελέσματα οδηγούν σε τέσσερα συνολικά συμπεράσματα.

\textbf{Πρώτον}, η παθητική κατάρρευση αποτελεί πραγματικό και επαναλαμβανόμενο φαινόμενο στο περιβάλλον κρυπτονομισμάτων, εκδηλούμενο τόσο σε αρχιτεκτονικές MLP όσο και σε LSTM όταν η συνάρτηση ανταμοιβής δεν παρέχει σαφές σήμα σχετικής απόδοσης. Η ανταμοιβή ενεργής απόδοσης αποτελεί την αποτελεσματικότερη αντιμετώπισή της.

\textbf{Δεύτερον}, η επιλογή της διατύπωσης του Λαγκρανζιανού περιορισμού είναι κρίσιμη. Ακατάλληλες διατυπώσεις (EMA, ομαλή κυρτή έλλειψη) μπορούν να δημιουργήσουν νέες μορφές κατάρρευσης που εμφανίζονται ως επιτυχία στις τυπικές μετρικές. Μόνο η διαγνωστική ανάλυση συμπεριφοράς (κατανομές ενεργειών, εξέλιξη μετρητών) αποκαλύπτει τη φύση της κατάρρευσης. Η soft-start διατύπωση, με βελτιστοποίηση μέσω CPCV, αποφεύγει αυτές τις παθολογίες.

\textbf{Τρίτον}, το σήμα συναισθήματος FinGPT παρέχει σημαντική αξία μέσω αλληλεπίδρασης με τον Λαγκρανζιανό περιορισμό liquify, και όχι ως αυτόνομο σήμα. Η συνέργεια (πίεση επένδυσης από τον περιορισμό ΚΑΙ κατεύθυνση επένδυσης από το συναίσθημα) αποφέρει δραματική βελτίωση Sharpe (+0,759 έναντι της ίδιας διαμόρφωσης χωρίς συναίσθημα), ενώ η απομονωμένη χρήση του σήματος αποφέρει οριακή μόνο βελτίωση (+0,030).

\textbf{Τέταρτον}, η βελτιστοποίηση μέσω CPCV/Optuna είναι αναγκαία για να αναδειχθεί η πλήρης αξία κάθε διαμόρφωσης: ο soft-start buffer χωρίς βελτιστοποίηση αποδίδει Sharpe 0,693, ενώ με CPCV βελτιστοποίηση φτάνει 1,026. Το CPCV εξασφαλίζει ότι η αναζήτηση δεν οδηγεί σε υπερπροσαρμογή στην περίοδο εκπαίδευσης, παρά το μεγάλο χώρο υπερπαραμέτρων.

Το κεντρικό αποτέλεσμα της εργασίας είναι ο πράκτορας Lagrangian liquify 15\%/70\% με χαρακτηριστικά FinGPT one-hot συναισθήματος, βελτιστοποιημένος μέσω CPCV/Optuna. Ο πράκτορας επιτυγχάνει Sharpe ratio 1,480 και τελική αξία χαρτοφυλακίου 1.496K από αρχικό κεφάλαιο 1.000K, υπερβαίνοντας κατά πολύ τόσο την ισόποση αγορά-και-κράτηση (922K, Sharpe 0,267) όσο και τον πράκτορα ενεργής απόδοσης χωρίς περιορισμούς ή συναίσθημα (1.335K, Sharpe 0,778). Αυτή η βελτίωση είναι αποτέλεσμα της συνδυαστικής εφαρμογής και των τριών αξόνων της εργασίας και δεν μπορεί να αποδοθεί σε κανέναν μεμονωμένα.
